% p211_1; source PDF page 211; manually verified editable table.
\begin{tikzpicture}[x=1cm,y=-1cm,inner sep=0pt,outer sep=0pt]
\path[use as bounding box] (0,0) rectangle (16.000,7.100);
\draw[line width=.4pt] (0.000,0.000) rectangle (5.200,1.200);
\node[align=center,text width=4.940cm,font={\linespread{1}\fontsize{10}{13.50}\selectfont\bfseries}] at (2.600,0.600) {};
\draw[line width=.4pt] (5.200,0.000) rectangle (10.400,1.200);
\node[align=center,text width=4.940cm,font={\linespread{1}\fontsize{10}{13.50}\selectfont\bfseries}] at (7.800,0.600) {缓刑};
\draw[line width=.4pt] (10.400,0.000) rectangle (16.000,1.200);
\node[align=center,text width=5.340cm,font={\linespread{1}\fontsize{10}{13.50}\selectfont\bfseries}] at (13.200,0.600) {假释};
\draw[line width=.4pt] (0.000,1.200) rectangle (5.200,2.550);
\node[align=left,text width=4.940cm,font={\linespread{1}\fontsize{10}{13.50}\selectfont}] at (2.600,1.875) {考验期内发现漏罪};
\draw[line width=.4pt] (5.200,1.200) rectangle (10.400,2.550);
\node[align=left,text width=4.940cm,font={\linespread{1}\fontsize{10}{13.50}\selectfont}] at (7.800,1.875) {撤销缓刑，数罪并罚，但不存在“减”去已执行刑期的问题};
\draw[line width=.4pt] (10.400,1.200) rectangle (16.000,2.550);
\node[align=left,text width=5.340cm,font={\linespread{1}\fontsize{10}{13.50}\selectfont}] at (13.200,1.875) {撤销假释，数罪并罚，先并后减};
\draw[line width=.4pt] (0.000,2.550) rectangle (5.200,3.400);
\node[align=left,text width=4.940cm,font={\linespread{1}\fontsize{10}{13.50}\selectfont}] at (2.600,2.975) {考验期满后发现漏罪};
\draw[line width=.4pt] (5.200,2.550) rectangle (10.400,3.400);
\node[align=left,text width=4.940cm,font={\linespread{1}\fontsize{10}{13.50}\selectfont}] at (7.800,2.975) {不能撤销，另行起诉};
\draw[line width=.4pt] (10.400,2.550) rectangle (16.000,3.400);
\node[align=left,text width=5.340cm,font={\linespread{1}\fontsize{10}{13.50}\selectfont}] at (13.200,2.975) {不能撤销，另行起诉};
\draw[line width=.4pt] (0.000,3.400) rectangle (5.200,4.900);
\node[align=left,text width=4.940cm,font={\linespread{1}\fontsize{10}{13.50}\selectfont}] at (2.600,4.150) {考验期内又犯新罪，考验期内或考验期满后发现};
\draw[line width=.4pt] (5.200,3.400) rectangle (10.400,4.900);
\node[align=left,text width=4.940cm,font={\linespread{1}\fontsize{10}{13.50}\selectfont}] at (7.800,4.150) {撤销缓刑，数罪并罚，但不存在“减”去已执行刑期的问题};
\draw[line width=.4pt] (10.400,3.400) rectangle (16.000,4.900);
\node[align=left,text width=5.340cm,font={\linespread{1}\fontsize{10}{13.50}\selectfont}] at (13.200,4.150) {撤销假释，数罪并罚，先减后并};
\draw[line width=.4pt] (0.000,4.900) rectangle (5.200,5.750);
\node[align=left,text width=4.940cm,font={\linespread{1}\fontsize{10}{13.50}\selectfont}] at (2.600,5.325) {考验期满后又犯新罪};
\draw[line width=.4pt] (5.200,4.900) rectangle (10.400,5.750);
\node[align=left,text width=4.940cm,font={\linespread{1}\fontsize{10}{13.50}\selectfont}] at (7.800,5.325) {不能撤销，也不可能构成累犯};
\draw[line width=.4pt] (10.400,4.900) rectangle (16.000,5.750);
\node[align=left,text width=5.340cm,font={\linespread{1}\fontsize{10}{13.50}\selectfont}] at (13.200,5.325) {不能撤销，但可能构成累犯};
\draw[line width=.4pt] (0.000,5.750) rectangle (5.200,7.100);
\node[align=left,text width=4.940cm,font={\linespread{1}\fontsize{10}{13.50}\selectfont}] at (2.600,6.425) {考验期内违反规定，考验期内或考验期满后发现};
\draw[line width=.4pt] (5.200,5.750) rectangle (10.400,7.100);
\node[align=left,text width=4.940cm,font={\linespread{1}\fontsize{10}{13.50}\selectfont}] at (7.800,6.425) {撤销缓刑};
\draw[line width=.4pt] (10.400,5.750) rectangle (16.000,7.100);
\node[align=left,text width=5.340cm,font={\linespread{1}\fontsize{10}{13.50}\selectfont}] at (13.200,6.425) {撤销假释};
\draw[line width=.4pt] (0,0)--(5.2,1.2);\node[font=\bfseries\fontsize{10}{13.5}\selectfont] at (3.8,.3) {处理方法};\node[font=\bfseries\fontsize{10}{13.5}\selectfont] at (1.8,.88) {期间情形};
\end{tikzpicture}
